%%%%%%%%%%%%%%%%%%%%%%%%%%%%%%%%%%%%%%%%%
% "ModernCV" CV and Cover Letter
% LaTeX Template
% Version 1.3 (29/10/16)
%
% This template has been downloaded from:
% http://www.LaTeXTemplates.com
%
% Original author:
% Xavier Danaux (xdanaux@gmail.com) with modifications by:
% Vel (vel@latextemplates.com)
%
% License:
% CC BY-NC-SA 3.0 (http://creativecommons.org/licenses/by-nc-sa/3.0/)
%
% Important note:
% This template requires the moderncv.cls and .sty files to be in the same 
% directory as this .tex file. These files provide the resume style and themes 
% used for structuring the document.
%
%%%%%%%%%%%%%%%%%%%%%%%%%%%%%%%%%%%%%%%%%

%----------------------------------------------------------------------------------------
%	PACKAGES AND OTHER DOCUMENT CONFIGURATIONS
%----------------------------------------------------------------------------------------

\documentclass[11pt,a4paper,sans]{moderncv} % Font sizes: 10, 11, or 12; paper sizes: a4paper, letterpaper, a5paper, legalpaper, executivepaper or landscape; font families: sans or roman

\moderncvstyle{classic} % CV theme - options include: 'casual' (default), 'classic', 'oldstyle' and 'banking'
\moderncvcolor{grey} % CV color - options include: 'blue' (default), 'orange', 'green', 'red', 'purple', 'grey' and 'black'

\usepackage{lipsum} % Used for inserting dummy 'Lorem ipsum' text into the template

\usepackage[scale=0.75]{geometry} % Reduce document margins
%\setlength{\hintscolumnwidth}{3cm} % Uncomment to change the width of the dates column
%\setlength{\makecvtitlenamewidth}{10cm} % For the 'classic' style, uncomment to adjust the width of the space allocated to your name

%----------------------------------------------------------------------------------------
%	NAME AND CONTACT INFORMATION SECTION
%----------------------------------------------------------------------------------------

\firstname{Tim} % Your first name
\familyname{Felgentreff} % Your last name

% All information in this block is optional, comment out any lines you don't need
\title{Curriculum Vit\ae{}}
\address{Straße zum Löwen 3}{14109 Berlin}
\mobile{0177 33 81 411}
% \phone{(000) 111 1112}
% \fax{(000) 111 1113}
\email{timfelgentreff@gmail.com}
\homepage{https://timfelgentreff.de}{timfelgentreff.de} % The first argument is the url for the clickable link, the second argument is the url displayed in the template - this allows special characters to be displayed such as the tilde in this example
% \extrainfo{}
\photo[70pt][0.4pt]{pictures/picture} % The first bracket is the picture height, the second is the thickness of the frame around the picture (0pt for no frame)
% \quote{"A witty and playful quotation" - John Smith}

%----------------------------------------------------------------------------------------

\begin{document}

%----------------------------------------------------------------------------------------
%	COVER LETTER
%----------------------------------------------------------------------------------------

% To remove the cover letter, comment out this entire block

\clearpage

\recipient{HR Department}{Corporation\\123 Pleasant Lane\\12345 City, State} % Letter recipient
\date{\today} % Letter date
\opening{Sehr geehrte Damen und Herren,} % Opening greeting
\closing{Mit freundlichen Grüßen,} % Closing phrase
\enclosure[Anhänge]{Curriculum Vit\ae{} (inkl. Lehrportfolio und Publikationsliste)\\%
  \hspace{4em} MSc Zeugnis\\
  \hspace{4em} Dissertationsurkunde\\
  \hspace{4em} Erklärung zu Lehr- und Forschungsvorhaben\\
  \hspace{4em} Publikationsproben} % List of enclosed documents

\makelettertitle % Print letter title

what then what then what then tail do I are you saying garnished garnished

Hi you get his voice
Dear Sir or Madam,
After my dissertation I joined Oracle labs has a senior member of the technical staff I joined Oracle apps for the research opportunities it provided and because it allowed me to focus on research topics such as virtual machine design programming tools damn and language design that I am interested in. Play during my time at Oracle apps I was the team lead for the Python implementation on top of the high performance grall VM compiler undo last three words and I was the principle investigator on the Oracle side for multiple research projects together with the hassle plattner Institute and the University of Cape Town ha ha yeah tio stop talking that the computer is recognizing your boys and it's saying you're saying yeah yeah haha Crash

Living history genius computing systems are here while we do provide programming models and tools to make the best use of these heterogeneous programming systems more and more non professional programmers with little or no formal education and programming are using heterogenous systems to process large amounts of data these users tend to focus on high level dynamic programming language for the ease of use hello this poses new challenges two optimizing compilers and runtime systems stop fat laptop XX laugh runtime optimizations have improved from localized peephole style optimizations on byte codes 2 complex heuristics involving branch duplication splitting and advanced inlining panjaa scape analysis techniques however all of these techniques that are in widespread use in the industry still focus on a small number off compute engines with future genius systems in widespread use optimizing runtime systems must provide the tools for the users to make optimum use of history genius hardware and automatically optimize code to run not only as fast as possible but also on the fastest available hardware. In order to enable these advanced optimizations hello laugh higher level abstractions must be optimized than previously that is the state of the art.

Give a little more programmers today have little or no formal programming education programming languages must provide high level must provide powerful high level abstractions in order to enable the users love to make use of advanced hardware capabilities such high level abstractions include declarative high specifications instead of imperative rules such as constraints and reactive programming systems Food

\lipsum[1-2] % Dummy text
\lipsum[4] % Dummy text

\makeletterclosing % Print letter signature

\newpage

%----------------------------------------------------------------------------------------
%	CURRICULUM VITAE
%----------------------------------------------------------------------------------------

\makecvtitle % Print the CV title

%----------------------------------------------------------------------------------------
%	EDUCATION SECTION
% ----------------------------------------------------------------------------------------

I am a member of the GraalVM team at Oracle Labs. I was led here by my interests
in programming languages and virtual machine design, and my involvement with the
PyPy/RPython project and various other dynamic language VMs. My work revolves
around dynamic language interoperability, connecting them with native C code in
a safe manner, and providing cross-language tool support. I also continue to
collaborate with the Hasso Plattner Institute at the University of Potsdam on
tools and environments for dynamic languagues and polyglot programming.

\section{Education}

\cventry{2012--2017}{Dr.rer-nat.}{Hasso Plattner Institute}{University of Potsdam}{Germany}{Summa Cum Laude}

\cvitem{Thesis}{The Design and Implementation of Object-Constraint Programming}

\cvitem{Advisor}{Robert Hirschfeld, Hasso Plattner Institute,
  University of Potsdam}

\cvitem{Reviewers}{Wolfgang de Meuter, Vrije Universiteit Brussel}
\cvitem{}{Alan Kay, Viewpoints Research Institute}

\cvitem{}{In my dissertation I presented Babelsberg, a formal design with
  multiple implementations for object-constraint programming to integrate
  constraint declaration and continuous satisfaction with mutable
  object-oriented structures and behavior.}

\cventry{2010--2012}{M.Sc.}{Hasso Plattner Institute}{University of
  Potsdam}{}{Grade: 1.0 (A+)}

\cvitem{Thesis}{Comparison, Replay, and Refinement of Communication Traces for
  Debugging Distributed Failures}

\cvitem{Advisor}{Robert Hirschfeld, Hasso Plattner Institute,
  University of Potsdam}

\cvitem{}{During my master's studies, I focused on programming language concepts
  to facilitate development tools, end-user programming, modularity, and
  meta-circularity especially in distributed and parallel programming settings.}

\cventry{2007--2010}{B.Sc.}{Hasso Plattner Institute}{University of
  Potsdam}{}{Grade: 1.4 (A)}

\cvitem{Thesis}{Post-mortem Analysis of Debug Traces}

\cventry{2005--2007}{Abitur}{Pückler Grammar School}{Cottbus}{Germany}{Grade: 1.0 (A+)}

\cventry{2004--2005}{AS-level studies}{Budehaven Community College}{Bude}{England}{}

\cventry{2000--2004}{GCSE-level studies}{Heinrich Heine Grammar School}{Cottbus}{Germany}{}

%----------------------------------------------------------------------------------------
%	WORK EXPERIENCE SECTION
%----------------------------------------------------------------------------------------

\section{Experience}

\cventry{July 2017--present}{Senior Software Engineer}{Oracle Labs}{Potsdam}{Germany}{}

\cvitem{}{At Oracle Labs I am a member of the Truffle/Graal group working on
  virtual machine technology for dynamic languages. I am the team lead for the
  high-performance Python implementation on GraalVM. I am also the Oracle PI for
  multiple research collaborations with the Hasso Plattner Institute on tools
  and environments for polyglot programming and with the University of Cape Town
  on native extension APIs for managed programming runtimes.}

\cventry{Apr 2016--June 2017}{Research Assistant}{Hasso Plattner
  Institute}{University of Potsdam}{Germany}{}

\cvitem{}{My research focused on novel programming paradigms, virtual machines
  for dynamic programming languages, and end-user programming with a special
  interest in computer science education for children.}

\cventry{Oct 2012--Apr 2016}{Ph.D. Student}{Hasso Plattner
  Institute}{University of Potsdam}{Germany}{}

\cvitem{}{My Ph.D. research focused on object-constraint programming, a novel
  approach at combining constraint solvers with dynamic, object-oriented
  languages in a clean and powerful way.}

\cventry{Oct 2014--Dec 2014}{Research Visit}{Nanjing University}{Nanjing}{China}{}

\cvitem{}{At Nanjing University, besides practicing Mandarin, I collaborated
  with researchers on emotion analysis in Twitter-like social networks.}

\cventry{Feb 2010--Oct 2012}{Working Student}{Finn GmbH}{Berlin}{Germany}{}

\cvitem{}{At Finn GmbH I was involved in plugin development for the Ruby on
  Rails based Redmine issue system}

\cventry{Oct 2009--Oct 2012}{Working Student}{Hasso Plattner
  Institute}{University of Potsdam}{Germany}{}

\cvitem{}{I was a tutor at the Software Architecture Group, assisting the staff
  with undergraduate lectures on software engineering and software development
  methodologies, using Squeak/Smalltalk as the educational language.}

\cventry{Aug 2011--Nov 2011}{Intern}{VMware}{Beaverton, OR}{USA}{}

\cvitem{}{At VMware I was working on the MagLev Ruby VM built on top of the
  distributed GemStone/S Smalltalk object database system.}

\cventry{May 2010--Aug 2010}{Ruby Summer of Code}{}{}{}{}

\cvitem{}{As part of the 2010 Ruby Summer of Code, I added support for Ruby C
  extensions to JRuby by emulating the CRuby extension interface using Java and
  JNI}


\subsection{Vocational}

\cventry{2012--Present}{1\textsuperscript{st} Year Analyst}{\textsc{Lehman Brothers}}{Los Angeles}{}{Developed spreadsheets for risk analysis on exotic derivatives on a wide array of commodities (ags, oils, precious and base metals), managed blotter and secondary trades on structured notes, liaised with Middle Office, Sales and Structuring for bookkeeping.
\newline{}\newline{}
Detailed achievements:
\begin{itemize}
\item Learned how to make amazing coffee
\item Finally determined the reason for \textsc{PC LOAD LETTER}:
\begin{itemize}
\item Paper jam
\item Software issues:
\begin{itemize}
\item Word not sending the correct data to printer
\item Windows trying to print in letter format
\end{itemize}
\item Coffee spilled inside printer
\end{itemize}
\item Broke the office record for number of kitten pictures in cubicle
\end{itemize}}

%------------------------------------------------

\cventry{2011--2012}{Summer Intern}{\textsc{Lehman Brothers}}{Los Angeles}{}{Rated "truly distinctive" for Analytical Skills and Teamwork.}

%------------------------------------------------

\subsection{Miscellaneous}

\cventry{2010--2011}{}{}{}{}{Spent some time finding myself. This was a courageous endeavour that didn't have a job title. It was quite important to my overall development though so I'm adding it to my CV. Also it explains the gap in my otherwise stellar CV.}

\cventry{2009--2010}{Computer Repair Specialist}{Buy More}{Burbank}{}{Worked in the Nerd Herd and helped to solve computer problems. Allowed me to become expert in all forms of martial arts and weaponry.}

%----------------------------------------------------------------------------------------
%	AWARDS SECTION
%----------------------------------------------------------------------------------------

\section{Awards}

\cvitem{2011}{School of Business Postgraduate Scholarship}
\cvitem{2010}{Top Achiever Award -- Commerce}

%----------------------------------------------------------------------------------------
%	COMPUTER SKILLS SECTION
%----------------------------------------------------------------------------------------

\section{Computer skills}

\cvitem{Basic}{\textsc{java}, Adobe Illustrator}
\cvitem{Intermediate}{\textsc{python}, \textsc{html}, \LaTeX, OpenOffice, Linux, Microsoft Windows}
\cvitem{Advanced}{Computer Hardware and Support}

%----------------------------------------------------------------------------------------
%	COMMUNICATION SKILLS SECTION
%----------------------------------------------------------------------------------------

\section{Communication Skills}

\cvitem{2010}{Oral Presentation at the California Business Conference}
\cvitem{2009}{Poster at the Annual Business Conference in Oregon}

%----------------------------------------------------------------------------------------
%	LANGUAGES SECTION
%----------------------------------------------------------------------------------------

\section{Languages}

\cvitemwithcomment{English}{Mothertongue}{}
\cvitemwithcomment{Spanish}{Intermediate}{Conversationally fluent}
\cvitemwithcomment{Dutch}{Basic}{Basic words and phrases only}

%----------------------------------------------------------------------------------------
%	INTERESTS SECTION
%----------------------------------------------------------------------------------------

\section{Interests}

\renewcommand{\listitemsymbol}{-~} % Changes the symbol used for lists

\cvlistdoubleitem{Piano}{Chess}
\cvlistdoubleitem{Cooking}{Dancing}
\cvlistitem{Running}

% ----------------------------------------------------------------------------------------

\clearpage

\section{Research Plan}

Heterogenous computing hardware has become more widespread and affordable in
recent years. At the same time, end-user programmers and data scientists with
little or no formal programming education are increasingly writing programs to
process large amounts of data. While hardware vendors and the research community
have demonstrated libraries [OpenMP] and languages [CUDA, OpenCL] to take
advantage of new hardware capabilities, high-level "scripting" languages such as
Python and JavaScript are still popular among the scientific computing
community. These languages are inherently single-threaded and not designed for
high-performance computing. Although popular libraries such as NumPy,
Tensorflow, or scikit-learn [1,2,3] work around this problem to an extent, in
practice the highly optimized code in these packages is interleaved with
high-level language code and thus inherently bottlenecked.

Millions of dollars by Google, Mozilla, Microsoft and others have been poured
into making JavaScript VMs fast. However, the need to balance the needs of the
browser environment in which JavaScript is often used with the needs for a
high-performance data processing language require trade-offs between warmup time
and peak performance. Furthermore, the optimizations applied by optimizing
virtual machines for JavaScript, Python, and other dynamic languages are still
focused on making single-threaded execution on the CPU as fast as possible.

As the typical size of problems for data scientists and end user programmers
grows and these groups have access to increasingly affordable GPU and FPGA
hardware, scripting languages and their optimizing virtual machines need to
improve to provide good performance. I propose to conduct research towards this
goal in three areas:
\begin{itemize}
\item I want to continue research upon the state of the art to identify and
  optimize high level patterns and control flow in languages to make use of
  auxiliary compute hardware. Work done on loop fusion and lazy evaluation
  [Weld], Columnar objects [Mattis], storage strategies [Gulenko], Megaguards
  [Qunaibit], and library-specific optimizations such as Numba or CuPy offer
  promising starting points for more runtime optimization techniques
\item Declarative programming techniques such as logic and constraint
  programming [Babelsberg, Squander, Kaplan] as well as reactive programming
  [Ramson, Salvaneschi] offer the chance for programmers to express intent. This
  can allow optimizing compilers to produce code more flexibly than for
  imperative programming techniques where the algorithm to find the solution is
  dictated in the code. I plan to continue to work on such programming language
  features to make languages more expressive and allow for better optimization
  at the same time.
\item Popular scripting languages do not allow all programming paradigms to
  exist equally. Additionally, libraries of interest to users or with specific
  features may not be available for all languages. The increasingly common
  notebook-style development environments can be a starting point for combining
  multiple different languages in a convenient way [Niephaus]. I plan to conduct
  further research into polyglot programming environments and the optimization
  across language boundaries [Wuerthinger]
\end{itemize}


\end{document}
%%% Local Variables:
%%% mode: latex
%%% TeX-master: t
%%% End:
